% ==============================================================
\section{Complexity Analysis}
\label{sec:complexity_appendix}
% ==============================================================

\begin{table}[h]
\centering
\caption{Computational complexity of the proposed method.
$N_P$: number of transactions containing pattern $P$,
$W_{\mathrm{sp}} = \prod_{k=1}^{d-1} (W_k + 1)$,
$\widetilde N_P = N_P W_{\mathrm{sp}}$,
$I_P$: total number of intervals in the candidate interval map.}
\label{tab:complexity}
\small
\begin{tabular}{lll}
\toprule
Step & Complexity \\
\midrule
Construction of $L_P$ & $O(N_P \cdot W_{\mathrm{sp}})$ \\
Sorting of $L_P$ & $O(\widetilde N_P \log \widetilde N_P)$ \\
Temporal interval & \multirow{2}{*}{$O\!\left(\sum_{\bm{s}_{\mathrm{sp}}}|L_P(\bm{s}_{\mathrm{sp}})| \log |L_P(\bm{s}_{\mathrm{sp}})|\right)$} \\
\quad computation & \\
Candidate interval & \multirow{2}{*}{$O(I_P)$} \\
\quad map intersection & \\
\bottomrule
\end{tabular}
\end{table}

The overall complexity faithful to the current implementation is
\[
  O\!\left(
    \sum_{P \in \mathcal{F}}
      \bigl(
        N_P W_{\mathrm{sp}}
        +
        \widetilde N_P \log \widetilde N_P
        +
        I_P
      \bigr)
  \right)
\]
where $\mathcal{F}$ is the set of candidate patterns actually evaluated.
The key factor governing practical complexity is $|\mathcal{F}|$,
which is controlled by Apriori pruning (Lemma~\ref{lem:apriori}):
any pattern whose $(k{-}1)$-subsets have empty candidate interval maps is pruned
without evaluation, bounding $|\mathcal{F}|$ to the set of dense patterns.
In the worst case, when all itemsets are dense,
$|\mathcal{F}|$ grows exponentially in $|\mathcal{I}|$;
this is inherent to any exact enumeration algorithm and is
mitigated in practice by the sparsity of real transaction data.

% ==============================================================
\section{Proofs of Lemmas}
\label{sec:proofs}
% ==============================================================

\begin{proof}[Proof of Lemma~\ref{lem:apriori}]
For any transaction $n$, if $P \subseteq d_n$ then $Q \subseteq d_n$.
Therefore, for any fixed spatial start vector $\bm{s}_{\mathrm{sp}}$,
each element of $L_P(\bm{s}_{\mathrm{sp}})$
is also contained in $L_Q(\bm{s}_{\mathrm{sp}})$.
Hence, for any $s_0$,
$S_P(s_0,\bm{s}_{\mathrm{sp}}) \leq S_Q(s_0,\bm{s}_{\mathrm{sp}})$ holds.
\end{proof}

\begin{proof}[Proof of Lemma~\ref{lem:stride}]
Consider the rightmost $\theta$ events within the current window
$[l, l+W_0]$.
The leftmost of these is $L[l_{\mathrm{lo}} + (c-\theta)] = l_{\mathrm{dense}}$.
For any $l' \leq l_{\mathrm{dense}}$,
all $\theta$ events are $\geq l'$,
and since $l' \geq l$, we have $l'+W_0 \geq l+W_0$.
Therefore, these events do not leave the current window $[l,l+W_0]$
and are all contained in $[l', l'+W_0]$ as well.
Hence, at least $\theta$ events remain within the window at $l'$.
\end{proof}


% ==============================================================
\section{EX-1 Synthetic Data Generation Details}
\label{sec:ex1_data_details}
% ==============================================================

\paragraph{Background data generation.}
One transaction is placed at every grid point of a
$T{=}100 \times S{=}100$ grid.
With $|\mathcal{I}|{=}40$ items,
each item is independently sampled from
Bernoulli$(p_{\mathrm{item}}{=}0.15)$ per cell to form the background.
Cells where all items are absent are assigned item~1 by default.

\paragraph{Pattern embedding.}
GT patterns $P_k = \{2k{-}1,\, 2k\}$ ($k=1,\ldots,20$)
are assigned to spatial bands
$s \in [5(k{-}1),\; 5(k{-}1){+}4]$ (width 5 cells).
Within cells belonging to the embedding region of each shape,
item pair $P_k$ is co-embedded independently with
Bernoulli$(p_{\mathrm{embed}}{=}0.25)$.

\paragraph{Embedding shapes.}
Each shape specifies the spatio-temporal distribution of pattern $P_k$
within its spatial band
($b_s = 5(k-1)$, $b_e = b_s + 4$).

\begin{itemize}
\item \textbf{Shape A (rectangular)}:
  Uniformly embedded in all cells with
  $t \in [30,70]$ and $s \in [b_s, b_e]$.
  The simplest fixed rectangular region in both time and space.
\item \textbf{Shape B (drift)}:
  At each time step $t \in [30,70]$,
  a band of width 5 is centered at
  $s_c = \lfloor b_s + (t - 30) / 40 \cdot 2 \rceil$
  (linearly shifting from $b_s$ to $b_s{+}2$).
  The dense band drifts smoothly along the spatial axis over time.
\item \textbf{Shape C (intermittent)}:
  $t \in [15,35] \cup [65,85]$, $s \in [b_s, b_e]$.
  Dense activity occurs in the same spatial band
  during two independent time intervals.
\item \textbf{Shape D (expanding)}:
  $t \in [30,40]$: center cell $b_s{+}2$ only (width 1);
  $t \in [41,55]$: $s \in [b_s{+}1, b_e{-}1]$ (width 3);
  $t \in [56,70]$: $s \in [b_s, b_e]$ (width 5).
  The spatial extent grows stepwise as time progresses.
\end{itemize}

\paragraph{Evaluation setup.}
Three random seeds (seed $\in \{42,55,68\}$) are used per shape,
giving $4 \times 3 = 12$ datasets in total.
The proposed method uses $W = (20, 5)$ and $\theta = 30$
uniformly across all shapes.
Grid search parameter ranges are shown in Table~\ref{tab:params_ex3}.

\begin{table}[h]
\centering
\caption{Grid search parameter ranges for EX-1 baseline methods.}
\label{tab:params_ex3}
\small
\begin{tabular}{lp{0.62\columnwidth}}
\toprule
Method & Grid search targets \\
\midrule
Colocation   & PI threshold $\tau \in \{0.05, 0.07, \ldots, 0.49\}$ (23 values) \\
STCP-Miner   & PI threshold (same) $\times$ window width $W_t \in \{4, 13, 26, 54\}$ \\
MZoom        & Top block count $k \in \{1, 2, 3, 5\}$ \\
ST-DBSCAN    & $\varepsilon_1 \in \{1,2,3\}$, $\varepsilon_2 \in \{3,7,15\}$, min\_pts $\in \{2,5\}$ \\
\bottomrule
\end{tabular}
\end{table}

% ==============================================================
\section{EX-2 Synthetic Data Generation Details}
\label{sec:ex2_data_details}
% ==============================================================

In EX-2, synthetic data is generated by placing
one transaction at \textbf{every} grid point of a $d = 2$ grid
$\{0,\ldots,T{-}1\} \times \{0,\ldots,S{-}1\}$.
For each transaction,
each item is independently sampled from
Bernoulli$(p_{\mathrm{item}})$ with
vocabulary size $|\mathcal{I}| = 20$ and
background occurrence probability $p_{\mathrm{item}} = 0.10$.
Cells where all items are absent
are assigned item~1 by default,
guaranteeing exactly one transaction per cell.
This ensures the number of transactions is exactly $N = T \times S$.

No embedded patterns are injected.
Dense patterns arise naturally from
statistical fluctuations of the Bernoulli background.
Preliminary calibration showed that
a sufficient number of patterns is detected
at threshold $\theta = 15$ across all configurations.


% ==============================================================
\section{Real Dataset Details}
\label{sec:real_data_qualitative}\label{sec:dataset_details}
% ==============================================================

\paragraph{Dataset overview.}
The four real-world datasets used in EX-3 (\S\ref{sec:ex3}) are detailed below.

\begin{itemize}
\item \textbf{Chicago Crime}~\cite{chicagocrime}:
  Crime report data published by the City of Chicago.
  Includes date/time, latitude, longitude,
  and crime type (33 types).
  Time discretized into weekly bins;
  latitude and longitude into $0.02^\circ$ equal-width bins
  (grid size $261 \times 20 \times 21$).
\item \textbf{NYPD Crime}~\cite{nypdcrime}:
  Crime report data published by the New York City Police Department.
  Includes crime type (69 types).
  Weekly bins, $0.02^\circ$ spatial bins
  (grid size $261 \times 43 \times 57$).
\item \textbf{NYC 311}~\cite{nyc311}:
  Citizen complaint (311) data from New York City.
  Includes complaint type (232 types).
  Weekly bins, $0.02^\circ$ spatial bins
  (grid size $261 \times 43 \times 57$).
\item \textbf{USGS Earthquake}~\cite{usgseq}:
  Earthquake catalog published by the United States Geological Survey.
  Magnitude discretized into 5 bands and used as items.
  This is the only 4-dimensional dataset
  (time $\times$ latitude $\times$ longitude $\times$ depth),
  with grid size $1{,}096 \times 20 \times 20 \times 4$.
\end{itemize}

\begin{table}[h]
\centering
\caption{Experimental datasets.
$N$: number of transactions (non-empty cells),
\#Items: number of item types,
avg$\ell$: average items per transaction,
$d$: dimensionality.}
\label{tab:datasets}
\small
\begin{tabular}{lrrrr}
\toprule
Dataset & $N$ & \#Items & avg$\ell$ & $d$ \\
\midrule
Chicago Crime~\cite{chicagocrime}  &  47,217 &  33 &  7.59 & 3 \\
NYPD Crime~\cite{nypdcrime}        & 189,481 &  69 &  5.67 & 3 \\
NYC 311~\cite{nyc311}              & 135,490 & 232 & 18.14 & 3 \\
USGS Earthquake~\cite{usgseq}      &  13,151 &   5 &  1.17 & 4 \\
\bottomrule
\end{tabular}
\end{table}

\paragraph{Experimental parameters.}
Maximum pattern length $|P| \leq 3$ was used for all datasets.
The three 3-dimensional datasets (Chicago/NYPD/NYC~311) share
temporal window $W_0 = 4$ weeks and spatial windows $W_k = 3$ bins
($k = 1, 2$), with only threshold $\theta$ adjusted per dataset.
USGS Earthquake used $W_0 = 13$ weeks, $W_k = 2$ bins ($k = 1,2,3$),
and $\theta = 2$.

\begin{table}[h]
\centering
\caption{Detection results of the proposed method on real datasets.
  $\theta$: density threshold,
  Total: number of patterns (all sizes),
  time: execution time (seconds).}
\label{tab:realdata_npatterns}
\small
\begin{tabular}{lrrrr}
\toprule
Dataset & $W$ & $\theta$ & Total & time (s) \\
\midrule
Chicago Crime  & $(4,3,3)$    &   3 & 2,835 &   9.7 \\
NYPD Crime     & $(4,3,3)$    &   5 & 4,478 &  38.2 \\
NYC 311        & $(4,3,3)$    &  35 & 5,130 & 108.7 \\
USGS Earthq.   & $(13,2,2,2)$ &   2 &    18 &   0.5 \\
\bottomrule
\end{tabular}
\end{table}
